\chapter{Testing and Validation}

\section*{Introduction}
\addcontentsline{toc}{section}{Introduction}

This chapter summarizes the testing and validation activities performed for
Football OS.

\section{Testing Strategy}

The testing approach combined complementary levels: smoke checks for service
health, functional and integration tests for core workflows across Gateway,
Governance, Workspace, and infrastructure dependencies, and security checks for
protected routes and role restrictions, completed by focused manual UI
validation of principal frontend flows.

\section{Test Environment}

Table~\ref{tab:test-environment} presents the integrated validation
environment used during testing.

\begin{table}[!htbp]
\centering
\caption{Components involved in the test environment}
\label{tab:test-environment}
{
\footnotesize
\renewcommand{\arraystretch}{1.08}
\setlength{\tabcolsep}{3.6pt}
\begin{tabularx}{\textwidth}{|L{4.0cm}|Y|}
\hline
\textbf{Component} & \textbf{Role in testing} \\
\hline
Docker Compose & Starts local services and dependencies. \\
\hline
API Gateway & Routing and protected entry-point checks. \\
\hline
Governance Service & Roles, permissions, canonical data, and policy-call checks. \\
\hline
Workspace Service & Documents, events, notifications, search, and profile checks. \\
\hline
PostgreSQL & Persistence checks for governance and audit data. \\
\hline
MongoDB & Persistence checks for workspace aggregates and metadata. \\
\hline
MinIO & Binary document storage and retrieval checks. \\
\hline
Keycloak & Authentication and identity-context checks. \\
\hline
OPA & Policy decision evaluation for authorization cases. \\
\hline
Kafka & Validation of asynchronous event/signal behavior. \\
\hline
OnlyOffice & Online document opening/editing callback checks. \\
\hline
Angular frontend/browser & Manual user workflow and interface checks. \\
\hline
\end{tabularx}
}
\end{table}

\section{Test Cases}

Table~\ref{tab:test-cases} summarizes representative scenarios and their
current validation status.

\begin{table}[!htbp]
\centering
\caption{Representative test cases and current validation status}
\label{tab:test-cases}
{
\scriptsize
\renewcommand{\arraystretch}{1.05}
\setlength{\tabcolsep}{3.4pt}
\begin{tabularx}{\textwidth}{|C{0.7cm}|Y|L{1.4cm}|Y|C{1.2cm}|}
\hline
\textbf{ID} & \textbf{Test Scenario} & \textbf{Type} & \textbf{Expected Result} & \textbf{Status} \\
\hline
T01 & Gateway health check & Smoke & Gateway is healthy. & Passed \\
\hline
T02 & Governance health check & Smoke & Governance is healthy. & Passed \\
\hline
T03 & Workspace health check & Smoke & Workspace is healthy. & Passed \\
\hline
T04 & User login with valid account & Functional & Session is created. & Passed \\
\hline
T05 & Unauthenticated protected access & Security & Access is denied. & Passed \\
\hline
T06 & Upload document & Integration & File and metadata stored. & Passed \\
\hline
T07 & Open document in OnlyOffice & Integration & Editor opens securely. & Partial \\
\hline
T08 & Archive document & Functional & Document is archived. & Passed \\
\hline
T09 & Create event as Head Coach & Functional & Event is persisted. & Passed \\
\hline
T10 & Event creation by unauthorized staff & Security & Policy denies creation. & Passed \\
\hline
T11 & Role-filtered player profile consultation & Integration & Permitted sections only. & Manual \\
\hline
T12 & Notification consultation & Functional & Notifications accessible. & Manual \\
\hline
T13 & Workspace search & Functional & Accessible results returned. & Partial \\
\hline
T14 & Medical document access control & Security & Medical access restricted. & Partial \\
\hline
\end{tabularx}
}
\end{table}

These tests were executed in a local Docker-based environment using health
checks, authenticated frontend workflows, and backend responses through the API
Gateway. The objective was integration and role-based-access validation in
development conditions, not production benchmarking.

\section{Performance Evaluation}

In addition to functional validation, a small set of latency observations was
collected in the local Docker-based environment. These measurements provide an
indicative view of integration behavior and relative operation cost. They are
not production benchmarks because they were executed on a local machine with
development configuration and a limited dataset.

\begin{table}[!htbp]
\centering
\caption{Local performance observations}
\label{tab:local-performance-observations}
{
\scriptsize
\renewcommand{\arraystretch}{1.05}
\setlength{\tabcolsep}{3.4pt}
\begin{tabularx}{\textwidth}{|Y|C{1.8cm}|Y|}
\hline
\textbf{Operation} & \textbf{Average Time} & \textbf{Observation} \\
\hline
Gateway health check & 35 ms & Confirms service availability. \\
\hline
Governance health check & 42 ms & Confirms service reachability. \\
\hline
Workspace health check & 48 ms & Confirms service reachability. \\
\hline
Event creation & 210 ms & Includes validation and MongoDB persistence. \\
\hline
Document metadata creation & 260 ms & Excludes large binary file transfer. \\
\hline
Document upload & 1.8 s & Depends on file size and MinIO transfer. \\
\hline
OnlyOffice config generation & 320 ms & Includes authorization and config generation. \\
\hline
Workspace search & 180 ms & Measured on a local sample dataset. \\
\hline
\end{tabularx}
}
\end{table}

Gateway and service-health operations remain below 50 ms in the local
environment, while persistence, authorization, storage, and editor operations
require more time. Document upload is the most expensive operation due to file
transfer and object-storage interaction; broader load testing remains necessary
before production deployment.

\section{Validation Summary}

Validation confirmed coherent integrated behavior across platform services, with
core workflows checked for authentication, document lifecycle, events,
notifications, profile consultation, and policy-based protection. Priority
improvements remain broader automated coverage, production monitoring readiness,
and deeper end-to-end validation.
